\section{VOO Framework Design}

\subsection{Type-Aware Fields and Constructors}

VOO provides expressive field declarations with type annotations (\texttt{double\_t},
\texttt{int\_t}, \texttt{string\_t}, \texttt{bool\_t}, \texttt{list\_t}, \texttt{dict\_t},
\texttt{obj\_t}) that serve documentation and constructor generation purposes without
runtime type enforcement. Fields support a \texttt{-static} modifier for class-level
storage and \texttt{public}/\texttt{private} visibility blocks. Every field has a default
value. VOO automatically generates three constructor variants: positional (\texttt{new}),
no-argument (\texttt{new()}), and named-argument (\texttt{new.args}). Custom constructors
can be declared for specialized initialization logic, and they are ordinary procedures
returning list values.

\begin{lstlisting}[language=Tcl, caption={VOO class declaration with multiple constructor variants}]
voo::class Person {
    public {
        string_t name "unknown"
        int_t age 0
        double_t salary 50000.0
    }

    method greet {} {
        return "Hello, I'm [get.name $this]"
    }
}

set p1 [Person::new "Alice" 30 75000.0]        ;# Positional
set p2 [Person::new()]                         ;# Defaults
set p3 [Person::new.args -name "Bob" -age 35]  ;# Named
\end{lstlisting}

\subsection{Accessors: Get, Set, Update}

For each field, VOO generates three accessor types. Getters receive the object by value
and return the field via \texttt{lindex}. Setters receive the variable name and modify
in-place via \texttt{upvar}/\texttt{lset}, providing copy-on-write safety. Updaters
temporarily detach a field into a local variable during a user-supplied script, preventing
COW propagation to the entire object during nested modifications---the
\texttt{try}/\texttt{finally} construct guarantees field reattachment even on exceptions.
Static fields use the \texttt{class.get.}/\texttt{class.set.} prefix convention.

\subsection{Methods and Inheritance}

Methods are procedures declared via \texttt{method} with optional modifiers:
none (\texttt{\$this} by value), \texttt{-static} (no \texttt{this}),
\texttt{-upvar} (\texttt{this} by reference), \texttt{-update \{fields\}} (fields
detached during body), and \texttt{-override} (validates parent method exists).
VOO supports single inheritance through \texttt{-extends}, where child classes inherit
parent fields with sequential indices and parent accessors are automatically available.
Parent methods can be imported explicitly via \texttt{importMethods}. Multiple inheritance
is excluded by design.

\subsection{Virtual Polymorphism}

VOO supports runtime polymorphic dispatch through the \texttt{-virtual} class flag and
\texttt{-virtual} method flag. A virtual class stores the concrete class namespace name
at index~0 of every instance---all other field indices shift up by one. This tag is an
interned \texttt{Tcl\_Obj*} shared by all instances of the same class; copy-on-write
ensures it is never duplicated on object copy, and its embedding as a literal at
class-definition time keeps constructor cost identical to non-virtual classes.

\begin{lstlisting}[language=Tcl, caption={Virtual polymorphic dispatch example}]
voo::class Shape -virtual {
    method area -virtual {} { return 0.0 }
}
voo::class Circle -extends Shape {
    public { double_t radius 1.0 }
    method area -override {} {
        return [expr {3.14159 * [get.radius $this] ** 2}]
    }
}
set s [Circle::new 5.0]
puts [Shape::area $s] ;# dispatches to Circle::area -> 78.54
\end{lstlisting}

For each \texttt{-virtual} method, VOO generates a \texttt{base.<name>} proc holding the
original body for direct parent calls, and makes the method itself a dispatcher that reads
index~0 and routes to the concrete class, falling back to \texttt{base.<name>} otherwise.
Child classes inherit virtual status automatically; \texttt{-override} methods are
auto-promoted to dispatchers, ensuring correct dispatch through deep inheritance chains.

\subsection{Visibility}

Fields and methods in \texttt{private \{ \}} blocks receive a \texttt{my.}~prefix and are
not exported. This naming-convention approach provides documentation-level protection
consistent with Tcl's philosophy---Itcl's enforcement mechanisms were bypassable through
namespace manipulation anyway.
